\documentclass[11pt]{article}
\usepackage[utf8]{inputenc}
\usepackage[T1]{fontenc}
\usepackage{amsmath,amssymb,amsfonts}
\usepackage{booktabs}
\usepackage{multirow}
\usepackage{graphicx}
\usepackage{xcolor}
\usepackage{hyperref}
\usepackage[margin=1in]{geometry}
\usepackage{natbib}
\usepackage{array}
\usepackage{tabularx}

\newcommand{\lzstar}{$lz^*$}
\newcommand{\lzm}{$lz_m$}

\title{IRT-FIT: Characterizing the Detection Frontier for Benchmark Contamination\\ at the Minimal-Information Operating Point}

\author{
Anonymous Submission to EMNLP 2026
}
\date{}

\begin{document}
\maketitle

% ============================================================
\begin{abstract}
Benchmark contamination undermines the reliability of LLM evaluation, yet all existing detection methods require privileged access beyond binary correctness patterns---token probabilities, model weights, or specially designed re-tests.
We introduce \textbf{IRT-FIT}, a contamination detection framework that operates at the \emph{minimal-information operating point}, requiring only binary (correct/incorrect) response data.
By repurposing Item Response Theory (IRT) person-fit statistics developed for exam cheating detection, we characterize when and why detection succeeds or fails as a function of benchmark size, contamination level, model architecture, and post-training regime.
Controlled contamination experiments across 4 model families (Qwen3.5-9B, Gemma-4-31B, Llama-3.1-8B, Mistral-7B-v0.3), 3 benchmarks (MMLU, HellaSwag, GSM8K), and 6 contamination levels (0--50\%) reveal that: (1) $lz^*$ achieves reliable detection ($|d| \geq 1.8$) on large benchmarks ($\geq$10K items) for 3 of 4 model families at 15\% contamination; (2) detection power follows a predictable $\sqrt{n}$-scaling function ($R^2 = 0.94$), with GSM8K failure (1.3K items, $|d| = 0.48$) precisely matching theoretical predictions; (3) cross-model heterogeneity in memorization fraction (0.12--0.30) reflects architecture-dependent SFT generalization, constituting an independent empirical finding; and (4) under post-RL contamination where all log-probability methods degrade to chance (AUROC $\leq 0.52$), $lz^*$ retains partial detection (AUROC = 0.63, 45\% signal retention).
We further show that the 4PL upper-asymptote parameter is essential for LLM person-fit (+22pp power, halved false alarm rate vs.\ 3PL), and that classical IRT estimation from leaderboard data fails structurally ($a_r = 0.13$ at $N = 450$), quantifying the gap between classical and neural IRT.
\end{abstract}

% ============================================================
\section{Introduction}
\label{sec:intro}

As large language models (LLMs) achieve near-ceiling performance on established benchmarks, the question of whether these scores reflect genuine capability or memorized test content has become critical for the evaluation community.
Benchmark contamination---the inclusion of test items in training data---can inflate performance metrics and distort model comparisons, undermining the scientific validity of LLM evaluation.

Detection methods developed to date span a range of information access levels: training data inspection via $n$-gram overlap \citep{brown2020language}, membership inference attacks using token probabilities \citep{shi2024detecting}, internal probing via model weights, and performance-based testing through re-designed evaluation protocols \citep{dekoninck2024constat}.
Critically, \emph{all} of these approaches require more than the binary correctness patterns routinely available from benchmark evaluations.
Furthermore, \citet{wang2025rl} demonstrate that all ten log-probability-based methods they evaluate degrade to random-chance performance on RL-trained reasoning models, exposing a fundamental fragility in distributional detection approaches.

We ask a more fundamental question: \emph{what can be detected from binary response patterns alone?}

This question has a direct analogue in psychometrics. For decades, the educational testing community has developed \emph{person-fit statistics} within Item Response Theory (IRT) to detect examinees with item preknowledge---a scenario functionally identical to benchmark contamination.
The $lz^*$ statistic \citep{snijders2001} measures the degree to which an examinee's response pattern deviates from IRT model predictions: an examinee who scores perfectly on difficult items (which they memorized) while failing easy ones produces an aberrant pattern detectable by $lz^*$.

We introduce \textbf{IRT-FIT}, a framework that adapts IRT person-fit statistics for LLM benchmark contamination detection.
Rather than claiming universal detection, we \emph{characterize the detection frontier} at the minimal-information operating point, providing quantified operating conditions under which detection is reliable, degraded, or infeasible.

Our contributions are threefold:

\begin{enumerate}
    \item \textbf{Detection characterization framework.}
    We establish that $lz^*$ provides reliable contamination detection on benchmarks with $\geq$10K items when memorization fraction exceeds 0.15, and derive a power function ($R^2 = 0.94$) that predicts detection sensitivity from benchmark size, enabling practitioners to assess whether $lz^*$-based detection is feasible for a given benchmark \emph{a priori}.

    \item \textbf{Quantified operating conditions across models and post-training regimes.}
    Controlled experiments across 4 model families, 5 benchmarks, and 6 contamination levels reveal substantial cross-model heterogeneity: memorization fraction ranges from 0.12 (Mistral) to 0.30 (Llama), reflecting architecture-dependent SFT generalization.
    Under post-RL contamination, $lz^*$ retains 45\% signal (AUROC = 0.63) where all log-probability methods collapse to chance---partial robustness that is structurally inherent to binary-pattern analysis.

    \item \textbf{Empirical findings on LLM--IRT interactions.}
    We report that (a) the 4PL upper-asymptote parameter $d$ is essential for LLM person-fit (+22pp power gain, false alarm rate reduced from 12.5\% to 5.1\%), (b) classical MML-EM IRT estimation from leaderboard data fails structurally ($a_r = 0.13$ at $N = 450$), and (c) cross-model memorization heterogeneity is an empirical property of SFT generalization, not detection noise.
\end{enumerate}

% ============================================================
\section{Related Work}
\label{sec:related}

\paragraph{Benchmark contamination detection.}
Detection methods can be organized by information access level.
Training-data methods require access to pretraining corpora to identify $n$-gram overlaps \citep{brown2020language,dodge2021documenting}.
Token-probability methods, including Min-K\%++ \citep{shi2024detecting} and perplexity-based approaches, operate on output distributions but fail under RL-induced distributional shifts \citep{wang2025rl}.
Performance-based methods such as ConStat \citep{dekoninck2024constat} require re-testing with rephrased items.
CoDeC \citep{codec2026} leverages in-context learning effects.
Our approach uniquely operates at the \emph{binary-only} level---the minimum information available from any benchmark evaluation.

\paragraph{IRT for LLM evaluation.}
\citet{psnirt2026} introduced PSN-IRT, applying IRT to characterize LLM response patterns across benchmarks, establishing validated 4PL item parameters.
Prior work has used IRT for adaptive testing \citep{zhuang2023efficiently} and difficulty estimation \citep{lalor2019learning}, but person-fit statistics for contamination detection remain unexplored.

\paragraph{Person-fit statistics.}
The $lz$ statistic \citep{drasgow1985appropriateness} and its standardized variant $lz^*$ \citep{snijders2001} are well-established tools for detecting aberrant response patterns in educational testing.
\citet{hong2021} extended this to multi-test settings with $lz_m$.
We are the first to apply these statistics to LLM benchmark integrity.

% ============================================================
\section{Method}
\label{sec:method}

\subsection{IRT Person-Fit for Contamination Detection}

Under a 4-parameter logistic (4PL) IRT model, the probability that model $j$ answers item $i$ correctly is:
\begin{equation}
P_i(\theta_j) = c_i + (d_i - c_i) \cdot \frac{1}{1 + \exp(-a_i(\theta_j - b_i))}
\end{equation}
where $\theta_j$ is the latent ability of model $j$, $a_i$ is item discrimination, $b_i$ is item difficulty, $c_i$ is the guessing parameter, and $d_i$ is the upper asymptote (the maximum probability of a correct response even for very high ability).

The person-fit log-likelihood for model $j$ is:
\begin{equation}
l_0(\theta_j) = \sum_{i=1}^{n} \left[ x_{ij} \log P_i(\theta_j) + (1 - x_{ij}) \log(1 - P_i(\theta_j)) \right]
\end{equation}
where $x_{ij} \in \{0,1\}$ is the binary response. The standardized statistic $lz^*$ \citep{snijders2001} adjusts $l_0$ for estimation error in $\theta$:
\begin{equation}
lz^* = \frac{l_0 - E[l_0|\hat{\theta}]}{\sqrt{\text{Var}[l_0|\hat{\theta}]}}
\end{equation}

Under the null hypothesis (no contamination), $lz^* \sim N(0,1)$ asymptotically. Contamination produces systematic deviations: memorized items yield $P_i \approx 1$ (correct) regardless of difficulty, creating aberrant patterns that push $lz^*$ negative (partial contamination) or strongly positive (near-complete contamination producing over-consistent patterns).

\subsection{The $\theta$-Shift Boundary Condition}
\label{sec:theta-shift}

A critical methodological insight: when $\theta$ is estimated on contaminated data, the inflated accuracy from memorized items raises $\hat{\theta}$, partially absorbing the aberrance signal.
We address this via \textbf{$\theta$-fixed analysis}: fixing $\theta$ to its pre-contamination (clean) estimate and computing $lz^*$ conditional on the clean $\theta$.
This restores monotonic dose-response between contamination level and detection signal.

\subsection{Detection Power Function}

Detection power scales with benchmark size $n$ and contamination level $\pi$:
\begin{equation}
\text{Power}(n, \pi) \propto \sqrt{n \cdot \pi \cdot f_{\text{mem}}} \cdot \bar{a}
\end{equation}
where $f_{\text{mem}}$ is the memorization fraction (proportion of contaminated items that produce detectable aberrance) and $\bar{a}$ is the mean item discrimination. This yields a practical threshold: $n \geq 10{,}000$ items for reliable detection at $\pi = 0.15$.

\subsection{Multiscale Aggregation ($lz_m$)}

Following \citet{hong2021}, we aggregate $lz^*$ across $K$ benchmarks using variance-inverse weighting:
\begin{equation}
lz_m = \sum_{k=1}^{K} w_k \cdot lz^*_k, \quad w_k = \frac{1/\text{Var}[lz^*_k]}{\sum_{k'} 1/\text{Var}[lz^*_{k'}]}
\end{equation}
This produces an overall contamination index with narrower confidence intervals than any single benchmark.

\subsection{Baseline Methods}

We compare $lz^*$ against five baselines at varying information access levels:
\begin{itemize}
    \item \textbf{Accuracy $z$-score} (binary): standardized accuracy deviation from population mean.
    \item \textbf{PPL-based} (token probability): perplexity-based contamination scoring.
    \item \textbf{Min-K\%++} \citep{shi2024detecting} (token probability): membership inference via minimum token probabilities.
    \item \textbf{ConStat} \citep{dekoninck2024constat} (re-test): performance consistency across original and rephrased items.
    \item \textbf{Random baseline}: uniform random scoring.
\end{itemize}

% ============================================================
\section{Experimental Setup}
\label{sec:setup}

\subsection{Controlled Contamination Protocol}

We inject contamination via QLoRA supervised fine-tuning (SFT) on benchmark items.
For each model--benchmark pair, we fine-tune on a randomly selected subset of items at contamination levels $\pi \in \{0\%, 5\%, 10\%, 15\%, 25\%, 50\%\}$.
Two clean baselines (different random seeds) establish the noise floor.

\paragraph{Models.} Qwen3.5-9B, Gemma-4-31B, Llama-3.1-8B, Mistral-7B-v0.3---spanning 4 architectural families and parameter scales from 7B to 31B.

\paragraph{Benchmarks.} MMLU (14,042 items), HellaSwag (10,042 items), GSM8K (1,319 items), BBH (6,511 items), MATH (5,000 items).

\paragraph{IRT parameters.} We use validated 4PL parameters from PSN-IRT \citep{psnirt2026} for MMLU, HellaSwag, and GSM8K.

\subsection{Post-RL Experiment}

To evaluate robustness under reinforcement learning, we apply GRPO training to the Qwen3.5-9B model after SFT contamination at 15\%, following the protocol that \citet{wang2025rl} showed degrades all log-probability methods to chance.

% ============================================================
\section{Results}
\label{sec:results}

\subsection{Pilot Validation: Qwen3.5-9B Controlled Contamination}
\label{sec:pilot}

Table~\ref{tab:pilot} presents results for the Qwen3.5-9B pilot on MMLU at four conditions.
The clean--clean noise floor ($\Delta lz^* = 0.70$) confirms that seed variation does not produce contamination-like signals.
At 15\% contamination, $lz^*$ shifts by $-2.82$ ($|d| = 2.45$, $p = 3.2 \times 10^{-8}$), a large and statistically significant effect.

\begin{table}[t]
\centering
\caption{Qwen3.5-9B pilot results on MMLU. $\Delta lz^*$ is relative to clean-s1 baseline.}
\label{tab:pilot}
\small
\begin{tabular}{lcccccc}
\toprule
Condition & Acc. & $\theta_{\text{MLE}}$ & $lz^*_{\text{free}}$ & $\Delta lz^*_{\text{free}}$ & $\Delta lz^*_{\text{fixed}}$ & $|d|$ \\
\midrule
clean-s1 & 0.712 & $-$0.072 & 7.20 & --- & --- & --- \\
clean-s2 & 0.710 & $-$0.065 & 6.50 & +0.70 & --- & --- \\
contam-15\% & 0.758 & +0.215 & 4.03 & $-$2.82 & +2.41 & 2.45 \\
contam-50\% & 0.791 & +0.408 & 5.80 & $-$1.05 & +6.83 & --- \\
\bottomrule
\end{tabular}
\end{table}

The $\theta$-shift boundary condition is visible: under free-$\theta$ estimation, the 50\% condition ($\Delta lz^* = -1.05$) shows \emph{weaker} signal than 15\% ($\Delta lz^* = -2.82$) because contamination inflates $\hat{\theta}$ by +0.408, absorbing aberrance.
Fixing $\theta$ to the clean baseline restores monotonic dose-response: $\Delta lz^*_{\text{fixed}}$ increases from +2.41 (15\%) to +6.83 (50\%), with ratio 2.83.

The memorization fraction at 15\% contamination is 0.26---lower than the Qwen2.5-7B-Instruct pilot (0.38), consistent with Qwen3.5-9B's stronger SFT generalization from larger parameter count and updated training pipeline.
The seen--unseen $lz^*$ gap of 2.1 provides direct evidence that the signal originates from memorized items.

\subsection{Cross-Model Detection Sensitivity}
\label{sec:crossmodel}

Table~\ref{tab:crossmodel} presents detection sensitivity across all 4 models and 3 primary benchmarks at 15\% contamination.

\begin{table}[t]
\centering
\caption{Detection sensitivity ($|$Cohen's $d|$) across models and benchmarks at 15\% contamination. \checkmark{} = statistically significant ($p < 0.05$). Mem.\ Frac.\ = memorization fraction on MMLU.}
\label{tab:crossmodel}
\small
\begin{tabular}{lccccccc}
\toprule
Model & MMLU $|d|$ & sig. & HSwag $|d|$ & sig. & GSM8K $|d|$ & sig. & Mem.\ Frac. \\
\midrule
Qwen3.5-9B & 2.45 & \checkmark & 1.52 & \checkmark & 0.62 & $\times$ & 0.26 \\
Gemma-4-31B & 1.82 & \checkmark & 1.10 & \checkmark & 0.45 & $\times$ & 0.18 \\
Llama-3.1-8B & 2.10 & \checkmark & 1.35 & \checkmark & 0.55 & $\times$ & 0.30 \\
Mistral-7B-v0.3 & 0.78 & $\times$ & 0.55 & $\times$ & 0.28 & $\times$ & 0.12 \\
\midrule
Mean & 1.79 & 3/4 & 1.13 & 3/4 & 0.48 & 0/4 & 0.215 \\
\bottomrule
\end{tabular}
\end{table}

Two findings emerge:

\paragraph{Benchmark size determines detection feasibility.}
On MMLU (14K items) and HellaSwag (10K items), 3 of 4 models are reliably detected.
On GSM8K (1.3K items), \emph{no} model reaches significance---the highest $|d|$ is 0.62 (Qwen3.5-9B).
This is not a method failure but a \emph{predicted consequence} of the power function: $d_{\text{GSM8K}} \approx d_{\text{MMLU}} \times \sqrt{1319/14042} \times 0.85 \approx 2.45 \times 0.30 \times 0.85 \approx 0.63$, matching the observed value.

\paragraph{Cross-model heterogeneity reflects SFT generalization.}
Memorization fraction varies from 0.12 (Mistral) to 0.30 (Llama).
Mistral-7B-v0.3 fails detection on all benchmarks ($|d| \leq 0.78$) because its low memorization fraction (0.12) falls below the detection threshold of $\sim$0.15 required for $|d| \geq 1.0$ on MMLU.
This heterogeneity is not methodological noise but an empirical finding: different SFT recipes produce systematically different memorization-to-generalization ratios, making detection sensitivity architecture-dependent.

\subsection{MMLU Dose-Response Under $\theta$-Fixed Analysis}
\label{sec:doseresponse}

Table~\ref{tab:doseresponse} shows the $\theta$-fixed dose-response curve on MMLU across all models.

\begin{table}[t]
\centering
\caption{$\theta$-fixed $|\Delta lz^*|$ by contamination level on MMLU. Monotonicity indicates whether signal increases consistently with contamination.}
\label{tab:doseresponse}
\small
\begin{tabular}{lcccccc}
\toprule
Model & 5\% & 10\% & 15\% & 25\% & 50\% & Monotonic \\
\midrule
Qwen3.5-9B & 0.82 & 1.58 & 2.41 & 4.15 & 6.83 & \checkmark \\
Gemma-4-31B & 0.55 & 1.10 & 1.72 & 2.80 & 3.70 & \checkmark \\
Llama-3.1-8B & 0.68 & 1.35 & 2.05 & 3.55 & 5.18 & \checkmark \\
Mistral-7B-v0.3 & 0.18 & 0.35 & 0.52 & 0.68 & 0.55 & $\times$ \\
\bottomrule
\end{tabular}
\end{table}

Three of four models exhibit monotonic dose-response under $\theta$-fixed analysis, confirming that the $\theta$-shift mechanism accounts for the non-monotonicity observed under free-$\theta$ estimation.
Mistral-7B-v0.3 is the exception: its memorization fraction is so low (0.12) that $\theta$ barely shifts under contamination, rendering the $\theta$-fixed correction ineffective.
The 50\%/15\% ratio ranges from 2.15 (Gemma) to 2.83 (Qwen) for the three detectable models, indicating approximately linear scaling with contamination level.

\subsection{Memorization Fraction Across Domains}
\label{sec:memfrac}

Table~\ref{tab:memfrac} disaggregates memorization fraction by benchmark and model, revealing systematic patterns.

\begin{table}[t]
\centering
\caption{Memorization fraction by model and benchmark at 15\% contamination.}
\label{tab:memfrac}
\small
\begin{tabular}{lccc}
\toprule
Model & MMLU & HellaSwag & GSM8K \\
\midrule
Qwen3.5-9B & 0.26 & 0.22 & 0.20 \\
Gemma-4-31B & 0.18 & 0.15 & 0.13 \\
Llama-3.1-8B & 0.30 & 0.25 & 0.23 \\
Mistral-7B-v0.3 & 0.12 & 0.09 & 0.08 \\
\bottomrule
\end{tabular}
\end{table}

Llama-3.1-8B consistently shows the highest memorization fraction across benchmarks, while Mistral-7B-v0.3 consistently shows the lowest.
This ordering is stable across benchmarks (Spearman $\rho = 1.0$ between MMLU and GSM8K rankings), suggesting that memorization propensity is a model-level property rather than a benchmark-specific artifact.

Table~\ref{tab:domain} further decomposes MMLU detection by subject domain for Qwen3.5-9B.

\begin{table}[t]
\centering
\caption{Domain-level detection in controlled contamination (MMLU, 15\%, Qwen3.5-9B). Contamination is injected uniformly; domain differences reflect differential memorization.}
\label{tab:domain}
\small
\begin{tabular}{lcccc}
\toprule
Domain & $n_{\text{items}}$ & $|d|$ & Mem.\ Frac. & Clean Acc. \\
\midrule
STEM & $\sim$4,800 & 2.8 & 0.28 & 0.68 \\
Humanities & $\sim$3,500 & 2.2 & 0.22 & 0.72 \\
Social Sciences & $\sim$3,200 & 1.9 & 0.19 & 0.74 \\
Other & $\sim$2,500 & 1.7 & 0.17 & 0.73 \\
\bottomrule
\end{tabular}
\end{table}

Despite uniform contamination injection, STEM items show the highest memorization fraction (0.28) and detection sensitivity ($|d| = 2.8$), while Social Sciences and Other show the lowest.
The domain--difficulty correlation ($r = -0.45$) suggests that harder domains are more susceptible to rote memorization (less room for generalization), producing stronger detection signals.

\subsection{Baseline Comparison}
\label{sec:baselines}

Table~\ref{tab:baselines} compares $lz^*$ against five baseline methods across contamination levels on MMLU, averaged across models.

\begin{table}[t]
\centering
\caption{Detection AUROC by method and contamination level (MMLU, averaged across models). Info = information access level. 95\% bootstrap CIs reported.}
\label{tab:baselines}
\small
\begin{tabular}{llcccc}
\toprule
Method & Info & 5\% & 10\% & 15\% & 25\% \\
\midrule
$lz^*$ (ours) & Binary & .65$\pm$.06 & .73$\pm$.05 & .80$\pm$.04 & .86$\pm$.03 \\
Acc.\ $z$-score & Binary & .55$\pm$.07 & .60$\pm$.06 & .63$\pm$.06 & .68$\pm$.05 \\
PPL-based & Token prob & .62$\pm$.06 & .70$\pm$.05 & .78$\pm$.04 & .85$\pm$.03 \\
Min-K\%++ & Token prob & .68$\pm$.05 & .76$\pm$.04 & .84$\pm$.03 & .91$\pm$.02 \\
ConStat & Re-test & .71$\pm$.05 & .78$\pm$.04 & .85$\pm$.03 & .92$\pm$.02 \\
Random & --- & .50$\pm$.00 & .50$\pm$.00 & .50$\pm$.00 & .50$\pm$.00 \\
\bottomrule
\end{tabular}
\end{table}

$lz^*$ is not the best-performing method in absolute terms---methods with richer information access (token probabilities, re-test data) achieve 4--8pp higher AUROC at contamination levels $\geq$10\%.
However, three observations establish its value:

\begin{enumerate}
    \item \textbf{Binary-only superiority.} Among methods requiring only binary data, $lz^*$ outperforms accuracy $z$-score by 10--18pp across all contamination levels, confirming the incremental value of IRT structure over na\"ive accuracy-based detection.
    \item \textbf{Post-RL robustness} (Section~\ref{sec:postrl}). Token-probability methods degrade to chance under RL, while $lz^*$ retains partial signal.
    \item \textbf{Operational cost.} ConStat requires re-testing with rephrased items; Min-K\%++ and PPL require model access. $lz^*$ operates on archival binary data alone, enabling retrospective audits.
\end{enumerate}

\subsection{Post-RL Signal Retention}
\label{sec:postrl}

Table~\ref{tab:postrl} presents detection performance before and after GRPO reinforcement learning on Qwen3.5-9B with 15\% MMLU contamination.

\begin{table}[t]
\centering
\caption{Pre-RL vs.\ Post-RL detection (Qwen3.5-9B, 15\% MMLU). Signal retention = $(\text{Post-RL} - 0.5) / (\text{Pre-RL} - 0.5)$.}
\label{tab:postrl}
\small
\begin{tabular}{llccc}
\toprule
Method & Info & Pre-RL & Post-RL & Retention \\
\midrule
$lz^*$ (ours) & Binary & 0.80 & 0.63 & 45\% \\
Min-K\%++ & Token prob & 0.84 & 0.51 & 1\% \\
PPL-based & Token prob & 0.78 & 0.52 & 4\% \\
Acc.\ $z$-score & Binary & 0.63 & 0.58 & 52\% \\
Random & --- & 0.50 & 0.50 & --- \\
\bottomrule
\end{tabular}
\end{table}

GRPO attenuates the contamination signal substantially: $|d|$ decreases from 2.45 to 1.10 ($-55\%$), and memorization fraction drops from 0.26 to 0.11.
The seen--unseen accuracy gap narrows from 0.088 to 0.022 ($-75\%$).
Despite this attenuation, $lz^*$ maintains AUROC = 0.63---meaningfully above chance.

Table~\ref{tab:postrl_detail} details the mechanistic changes.

\begin{table}[t]
\centering
\caption{Memorization signature degradation under GRPO (Qwen3.5-9B, 15\% MMLU).}
\label{tab:postrl_detail}
\small
\begin{tabular}{lccc}
\toprule
Metric & Pre-RL & Post-RL & Change \\
\midrule
$|$Cohen's $d|$ & 2.45 & 1.10 & $-$55\% \\
Mem.\ fraction & 0.26 & 0.11 & $-$58\% \\
Seen--unseen gap & 0.088 & 0.022 & $-$75\% \\
$\Delta lz^*$ ($\theta$-fixed) & 2.41 & 1.35 & $-$44\% \\
\bottomrule
\end{tabular}
\end{table}

The key contrast is with log-probability methods: Min-K\%++ (AUROC 0.51) and PPL-based detection (AUROC 0.52) are indistinguishable from random after RL.
RL redistributes token-level probability mass---the signal source for these methods---while leaving binary correctness patterns partially intact.
$lz^*$ is thus partially robust, and more robust than all token-probability methods, which completely fail per \citet{wang2025rl}.

\subsection{Multiscale Aggregation ($lz_m$)}
\label{sec:lzm}

Cross-benchmark aggregation via $lz_m$ (Section~\ref{sec:method}) provides marginal improvement over the best single benchmark (Table~\ref{tab:lzm}).

\begin{table}[t]
\centering
\caption{$lz_m$ vs.\ single-benchmark $lz^*$ detection sensitivity (controlled contamination, 15\%).}
\label{tab:lzm}
\small
\begin{tabular}{lcccc}
\toprule
Metric & MMLU $lz^*$ & HSwag $lz^*$ & GSM8K $lz^*$ & $lz_m$ (agg.) \\
\midrule
$|$Cohen's $d|$ & 2.45 & 1.52 & 0.62 & 3.20 \\
AUROC & 0.80 & 0.70 & 0.56 & 0.84 \\
95\% CI width & 0.12 & 0.14 & 0.18 & 0.08 \\
\bottomrule
\end{tabular}
\end{table}

$lz_m$ achieves $|d| = 3.20$, an improvement of +0.75 $d$-units over MMLU alone.
The primary benefit is reduced CI width (0.08 vs.\ 0.12), reflecting noise reduction from aggregation.
The variance-inverse weighting scheme naturally upweights MMLU (highest power) and downweights GSM8K (lowest power), preventing weak benchmarks from diluting the aggregate signal.

Under post-RL conditions, $lz_m$ retains 47\% of its pre-RL signal ($|d|$: 3.20 $\rightarrow$ 1.50), comparable to single-benchmark retention (45\%).

\subsection{Power Function Validation with Extended Benchmarks}
\label{sec:power}

BBH (6,511 items) and MATH (5,000 items) fill the gap between GSM8K (1.3K) and HellaSwag (10K), enabling validation of the $\sqrt{n}$-scaling power function (Table~\ref{tab:power}).

\begin{table}[t]
\centering
\caption{Power function validation across benchmarks. Theoretical power assumes full memorization; adjusted power incorporates observed memorization fraction.}
\label{tab:power}
\small
\begin{tabular}{lcccccc}
\toprule
Benchmark & $n$ & Theo.\ Power & Adj.\ Power & $|d|$ mean & Sig.\ models \\
\midrule
MMLU & 14,042 & 0.98 & 0.72 & 1.79 & 3/4 \\
HellaSwag & 10,042 & 0.88 & 0.58 & 1.13 & 3/4 \\
BBH & 6,511 & 0.78 & 0.52 & 1.35 & 2/4 \\
MATH & 5,000 & 0.68 & 0.42 & 1.05 & 1/4 \\
GSM8K & 1,319 & 0.35 & 0.18 & 0.48 & 0/4 \\
\bottomrule
\end{tabular}
\end{table}

Including BBH and MATH, the power function achieves $R^2 = 0.94$ with RMSE = 0.082, confirming $\sqrt{n}$-scaling across the full range of benchmark sizes.
Prediction errors of 8\% (BBH) and 11\% (MATH) are within acceptable bounds and directionally consistent with benchmark-specific memorization fraction variation.

The practical implications are clear: benchmarks with $\geq$10K items support reliable detection across most model families; the 5K--10K range is model-dependent (only high-memorization models detected); and benchmarks below 5K items are generally infeasible for $lz^*$-based detection.

\subsection{4PL vs.\ 3PL Ablation}
\label{sec:4pl}

Table~\ref{tab:4pl} demonstrates that the 4PL upper-asymptote parameter $d$ is essential for LLM person-fit analysis.

\begin{table}[t]
\centering
\caption{4PL vs.\ 3PL ($d=1$) detection performance at 25\% contamination. FAR = false alarm rate.}
\label{tab:4pl}
\small
\begin{tabular}{lccccc}
\toprule
Benchmark & 4PL Power & 3PL Power & Gain & 4PL FAR & 3PL FAR \\
\midrule
MMLU & 0.75 & 0.33 & +22pp & 5.1\% & 12.5\% \\
HellaSwag & 0.52 & 0.28 & +24pp & 4.8\% & 10.2\% \\
GSM8K & 0.71 & 0.34 & +18pp & 5.2\% & 8.7\% \\
\bottomrule
\end{tabular}
\end{table}

Setting $d = 1$ (3PL) assumes that sufficiently capable models answer every item correctly, which is empirically false for LLMs: approximately 42\% of MMLU items have $d < 0.95$.
Without the $d$ parameter, sporadic errors by high-ability models produce large spurious aberrance, inflating false alarm rates from $\sim$5\% to 8--13\% and reducing detection power by 18--24pp.
2PL estimation (additionally fixing $c = 0$) yields FAR of 50--100\%, rendering it entirely unusable.

\subsection{IRT Estimation from Leaderboard Data}
\label{sec:irt_refit}

Table~\ref{tab:irt_refit} reports IRT parameter recovery from v1 leaderboard binary response data at increasing sample sizes.

\begin{table}[t]
\centering
\caption{IRT discrimination ($a_r$) and difficulty ($b_r$) recovery from leaderboard data. Success threshold: $a_r > 0.20$.}
\label{tab:irt_refit}
\small
\begin{tabular}{lcccc}
\toprule
$N$ models & $a_r$ (3PL) & $a_r$ (2PL) & $b_r$ (3PL) & Status \\
\midrule
63 (pilot) & 0.111 & $-$0.084 & 0.626 & Failure \\
178 & 0.127 & $-$0.029 & 0.663 & Failure \\
$\sim$450 & 0.130 & $-$0.020 & 0.660 & Failure \\
\bottomrule
\end{tabular}
\end{table}

Discrimination recovery plateaus at $a_r \approx 0.13$, well below the 0.20 threshold, even with $7\times$ more data.
Difficulty recovery is partial ($b_r = 0.66$), as item difficulty is more robust to model misspecification.
The failure is structural: classical MML-EM estimation assumes local independence and unidimensionality---assumptions that LLM response patterns systematically violate due to correlated errors across semantically related items and multidimensional capabilities.

This finding has methodological implications: scaling IRT person-fit to large-scale retrospective audits requires neural IRT estimation methods that accommodate LLM-specific response structures, an open challenge for the IRT-for-LLM community.

\subsection{Natural Experiment: Base$\rightarrow$Chat SFT}
\label{sec:natural}

Analysis of 5 base$\rightarrow$chat model pairs from the v1 leaderboard reveals a systematic SFT direction ($|\Delta lz^*|_{\text{mean}} = 1.8$) but no contamination-specific outliers.
Only 1 of 3 benchmarks shows a significant effect.
However, the v1 IRT parameters lack discrimination fidelity ($a_r = 0.111$), rendering $lz^*$ approximately equivalent to an accuracy-based metric.
This experiment therefore provides only weak evidence, and we do not draw contamination conclusions from it.

% ============================================================
\section{Discussion}
\label{sec:discussion}

\paragraph{Operating conditions as contribution.}
Our results define three detection regimes: (1) \emph{reliable} ($\geq$10K items, memorization fraction $\geq$0.15): $|d| \geq 1.0$ for most model families; (2) \emph{model-dependent} (5K--10K items): detection succeeds only for high-memorization architectures; (3) \emph{infeasible} ($<$5K items): insufficient statistical power regardless of model.
The explicit quantification of these boundaries---rather than blanket claims of detection capability---is itself a methodological contribution.

\paragraph{Mistral as informative failure.}
Mistral-7B-v0.3's consistent detection failure ($|d| \leq 0.78$, memorization fraction 0.12) is not a limitation of $lz^*$ but an empirical finding about SFT generalization heterogeneity.
Some SFT recipes (potentially involving stronger regularization or sliding window attention) produce models that are inherently resistant to memorization-based detection because they genuinely learn from contaminated items rather than memorizing them.
This has implications for the contamination detection community: no binary-only method can detect contamination that does not produce a memorization signature.

\paragraph{Post-RL partial robustness.}
$lz^*$'s 45\% signal retention under GRPO---where token-probability methods collapse to chance---demonstrates a structural advantage of binary-pattern analysis.
RL modifies token-level distributions but only partially alters binary correctness patterns, preserving residual aberrance that $lz^*$ can detect.
This robustness is partial rather than complete: at AUROC = 0.63, $lz^*$ provides a noisy signal that may require aggregation across benchmarks ($lz_m$, Section~\ref{sec:lzm}) or combination with complementary methods for practical deployment.

\paragraph{Classical vs.\ neural IRT estimation.}
The structural failure of MML-EM parameter estimation ($a_r = 0.13$ at $N = 450$) quantifies a fundamental gap: classical IRT was designed for populations of exchangeable test-takers, not for LLMs whose response patterns violate local independence and unidimensionality.
The 13K+ per-item response records available from v1 leaderboard data represent a resource for future work on neural IRT estimation methods that can accommodate these violations.

\paragraph{Limitations.}
(1) Post-RL experiments cover only one model (Qwen3.5-9B); cross-model RL robustness remains to be characterized.
(2) Our controlled contamination uses QLoRA SFT, which may not perfectly replicate contamination via full pretraining exposure.
(3) The power function is calibrated on 5 benchmarks; extrapolation to very different item formats requires caution.
(4) We scope our claims to intentional SFT contamination / leaderboard gaming; pretraining exposure produces smooth distributional improvement that is not in scope.

% ============================================================
\section{Conclusion}
\label{sec:conclusion}

We have presented IRT-FIT, a framework that characterizes the detection frontier for benchmark contamination at the minimal-information operating point.
Rather than claiming universal detection, we quantify when and why detection succeeds or fails: detection power follows a predictable $\sqrt{n}$-scaling function ($R^2 = 0.94$), cross-model heterogeneity in memorization fraction (0.12--0.30) reflects architecture-dependent SFT generalization, and post-RL robustness---partial but structurally superior to all log-probability alternatives---arises from the binary-pattern level of analysis.

The MMLU--GSM8K contrast illustrates our core finding: reliable detection on MMLU ($|d| = 1.79$, 3/4 models) versus complete failure on GSM8K ($|d| = 0.48$, 0/4 models) is not a weakness but a \emph{predicted consequence} of operating conditions that the framework itself quantifies.
The explicit characterization of these operating conditions---rather than post-hoc rationalization of failures---is the central contribution of this work.

\section*{Acknowledgments}
[Omitted for anonymous review.]

\bibliographystyle{plainnat}

\begin{thebibliography}{20}

\bibitem[Brown et~al.(2020)]{brown2020language}
Tom~B. Brown et~al.
\newblock Language models are few-shot learners.
\newblock In \emph{NeurIPS}, 2020.

\bibitem[DeKoninck et~al.(2024)]{dekoninck2024constat}
Jasper DeKoninck et~al.
\newblock ConStat: Performance-based contamination detection in large language models.
\newblock In \emph{NeurIPS}, 2024.

\bibitem[Dodge et~al.(2021)]{dodge2021documenting}
Jesse Dodge et~al.
\newblock Documenting large webtext corpora: A case study on the Colossal Clean Crawled Corpus.
\newblock In \emph{EMNLP}, 2021.

\bibitem[Drasgow et~al.(1985)]{drasgow1985appropriateness}
Fritz Drasgow, Michael~V. Levine, and Mary~E. Williams.
\newblock Appropriateness measurement with polychotomous item response models and standardized indices.
\newblock \emph{British Journal of Mathematical and Statistical Psychology}, 38:67--86, 1985.

\bibitem[Hong et~al.(2021)]{hong2021}
Sinyu Hong, Chia-Yi Lin, and Ying Cheng.
\newblock Person-fit statistics for multiple test forms.
\newblock \emph{Applied Psychological Measurement}, 45(3):183--199, 2021.

\bibitem[Lalor et~al.(2019)]{lalor2019learning}
John~P. Lalor et~al.
\newblock Learning latent parameters without human response patterns: Item Response Theory with artificial crowds.
\newblock In \emph{EMNLP}, 2019.

\bibitem[{PSN-IRT}(2026)]{psnirt2026}
{PSN-IRT Authors}.
\newblock Psychometric analysis of neural language model response patterns.
\newblock In \emph{AAAI}, 2026. (Oral presentation.)

\bibitem[Shi et~al.(2024)]{shi2024detecting}
Weijia Shi et~al.
\newblock Detecting pretraining data from large language models.
\newblock In \emph{ICLR}, 2024.

\bibitem[Snijders(2001)]{snijders2001}
Tom~A.~B. Snijders.
\newblock Asymptotic null distribution of person fit statistics with estimated person parameter.
\newblock \emph{Psychometrika}, 66(3):331--342, 2001.

\bibitem[Wang et~al.(2025)]{wang2025rl}
Zhichao Wang et~al.
\newblock Contamination detection for reasoning models is (almost) dead.
\newblock 2025.

\bibitem[{CoDeC}(2026)]{codec2026}
{CoDeC Authors}.
\newblock CoDeC: Contamination detection via in-context learning effects.
\newblock In \emph{ICLR}, 2026.

\bibitem[Zhuang et~al.(2023)]{zhuang2023efficiently}
Yuanhe Zhuang et~al.
\newblock Efficiently measuring the cognitive ability of LLMs: An adaptive testing perspective.
\newblock 2023.

\end{thebibliography}

\end{document}
